% ============================================================
%  附录
% ============================================================

\section{支撑材料文件列表}

% 按实际压缩包内容填写文件名、用途和是否可运行。
% 若确实没有支撑材料，请在此明确说明“本论文没有支撑材料”。
\begin{center}
\begin{tabularx}{0.92\textwidth}{cJY}
  \topline
  文件名 & 内容说明 & 备注 \\
  \midline
  code/preprocess.py & 数据预处理脚本 & 可运行 \\
  data/processed.csv & 预处理后的中间数据 & 由脚本生成 \\
  \bottomline
\end{tabularx}
\end{center}

\section{主要代码清单}

% 推荐使用 \lstinputlisting 引用外部代码文件（便于版本管理）：
% \lstinputlisting[language=Python, caption={数据预处理脚本}]{./code/preprocess.py}

% 内联代码示例：
\begin{lstlisting}[language=Python, caption={示例：读取数据并计算均值}]
import pandas as pd

def load_data(path: str) -> pd.DataFrame:
    return pd.read_csv(path)

if __name__ == "__main__":
    df = load_data("../examples/dummy_data/sample_series.csv")
    print(df["value"].mean())
\end{lstlisting}

注：完整的代码文件应在提交时一并打包。

\section{补充图表}

% 在此放置正文中未包含的补充性图表。
